\documentclass[../main]{subfiles}
\begin{document}
% ============================================================
% Chapter 4: Influence of Initial Stress Anisotropy on
%            Liquefaction Resistance
% ============================================================
\chapter{Influence of Initial Stress Anisotropy on Liquefaction Resistance}
\label{ch:anisotropic}

% ----------------------------------------------------------------
\section{Introduction}
\label{sec:ch4_intro}

During earthquake loading, saturated granular soils can experience a rapid rise in pore pressure that drastically reduces their shear strength. This phenomenon, termed liquefaction, has caused severe damage to infrastructure and buildings \citep{IshiharaKoga1981,SeedIdriss1967}. A large body of laboratory work has investigated liquefaction resistance using cyclic undrained triaxial tests, examining the effects of cyclic stress ratio (CSR), relative density, and confining pressure \citep{Hyodo1991,SeedLee1966,Silver1976,Toki1986,Yoshimi1984}. However, vertically propagating shear waves in the ground induce gradually varying shear stress on soil elements, leading to continuous rotation of principal stress axes \citep{Arthur1980,IshiharaTowhata1983,IshiharaYasuda1975}, which cannot be fully represented in conventional triaxial loading.

Hollow cylinder apparatus (HCA) and hollow torsional shear tests address this limitation by allowing torsional loading with principal stress axis rotation. Early studies using HCA-type devices \citep[e.g.,][]{IshiharaYasuda1975,Tatsuoka1986,YamashitaToki1993} showed that liquefaction behavior can differ from triaxial-test results and depends on the loading mode and stress path. Subsequent research by \citet{Tatsuoka1986} compared triaxial and torsional tests on specimens prepared with varying methods, revealing inconsistencies in liquefaction resistance between the two testing approaches. Further investigations by \citet{YamashitaToki1993} utilized torsional and triaxial tests to refine constitutive models for liquefiable sands. These studies demonstrated that the chosen testing method---whether torsional or triaxial---influenced the measured liquefaction resistance. These observations motivate the use of HCA-type boundary conditions when investigating liquefaction mechanisms under more realistic cyclic loading paths.

The influence of the initial lateral-to-vertical effective stress ratio, \Kzero{}, on liquefaction resistance has been investigated extensively, yet the reported findings vary considerably depending on test conditions. Some experimental studies reported weak or negligible dependence of liquefaction strength on the initial anisotropic stress state \citep[e.g.,][]{IshiharaTakatsu1979,YamashitaToki1993}, whereas others found clear differences between isotropically consolidated (IC) and anisotropically consolidated (AC) specimens, with the trend depending on density \citep{GeorgiannouKonstadinou2014}. \citet{Vargas2020} observed that specimens under compressive consolidation exhibited higher liquefaction resistance than those consolidated isotropically. More recently, \citet{Zhao2024} conducted systematic torsional shear tests covering both compressional and extensional consolidation states and attributed the higher liquefaction resistance under compressive consolidation to smaller principal stress axis rotation during cyclic loading. While these findings establish the macroscopic trend convincingly, the particle-scale mechanisms through which stress anisotropy alters liquefaction resistance remain to be explored.

The discrete element method \citep[DEM;][]{Cundall1979} is well suited for this problem because it provides direct access to particle-scale quantities such as contact topology, force chains, and fabric anisotropy, while avoiding some specimen-to-specimen variability associated with laboratory preparation. Numerous DEM studies have examined liquefaction-related behavior in triaxial, simple-shear, or polyhedral specimen settings \citep[e.g.,][]{Huang2018,Yang2021,Jiang2021,Morimoto2021,WuYang2022,Xie2023,YangHuang2023,Zhang2023} and have identified important microscopic factors governing liquefaction resistance. Recent DEM investigations have also addressed \Kzero{} effects: \citet{YangTaiebat2024} found that both preparation protocols and \Kzero{} influence liquefaction resistance, with the difference between IC and AC states narrowing at higher densities; \citet{Otsubo2022} showed that specimens with stronger inherent anisotropy exhibit weaker stiffness in the minor principal direction and correspondingly lower liquefaction resistance. However, triaxial configurations do not incorporate principal stress axis rotation, while periodic-boundary and polyhedral specimen settings, though capable of applying general stress paths, lack direct counterparts in standard laboratory apparatus. Although DEM replication of HCA tests has precedents under drained loading \citep[e.g.,][]{Li2014,Liu2021}, extension to undrained cyclic loading in HCA geometry remains limited.

Using DEM to replicate HCA tests is relatively specialized, but still has precedents. \citet{Li2014} conducted DEM simulations of drained tests and investigated strain localization in HCA tests. \citet{Liu2021} conducted analysis of torsional shear tests under drained conditions and investigated the development of cracks at different principal stress rotation angles. \citet{Han2024jgs} first realized undrained cyclic torsional shear in a DEM-based HCA model by introducing a combined servo mechanism that simultaneously satisfies the HCA stress boundary conditions and the constant-volume constraint. However, their formulation assumed a constant specimen height throughout, precluding axial loading and height variation, which are commonly employed loading modes in HCA tests. \citet{Ma2024} subsequently generalized this framework to accommodate a broader range of loading conditions, including drained stress paths and combined axial--torsional loading. While these formulations successfully reproduced a variety of target stress paths, three aspects of the undrained servo problem deserve closer examination:
\begin{enumerate}
\item[(a)] Under undrained loading, the constant-volume constraint is a governing equation, not an auxiliary condition. It must be substituted into the kinematic system to eliminate one degree of freedom before the stress boundary conditions can form a closed system.
\item[(b)] The servo conditions should reproduce the laboratory total-stress boundary conditions. Because the pore water pressure $u$ is a measured response rather than a prescribed input, directly targeting individual effective stresses, as in previous formulations, is only an approximation. A valid formulation must express the servo conditions in a form that does not require $u$ as an input.
\item[(c)] Under the constant-volume constraint, the kinematic degrees of freedom (inner radius, outer radius, and height) are geometrically coupled, so movement of one wall inevitably affects the stresses on others. The servo coefficients must therefore account for this cross-coupling, rather than being derived independently from the contact stiffness at each wall.
\end{enumerate}
This chapter makes this closure explicit and derives closed-form servo coefficients at every timestep by analytically inverting the coupled stiffness matrix. This analytical exactness is especially important near liquefaction, where rapid loss of particle contacts causes large timestep-to-timestep stiffness changes and accumulated servo error can compromise boundary-condition control---an instability regime also noted by \citet{Ma2024}.

The investigation is organized in two stages. First, the macroscopic \Kzero{}-dependent liquefaction resistance trend is established for representative compression, isotropic, and extension states across multiple cyclic stress ratios, accompanied by energy and stiffness analyses. Second, the parameter space is expanded to two relative densities and five \Kzero{} values, and microscopic fabric analyses in the cylindrical coordinate system are used to separate the respective roles of relative density and stress anisotropy at the particle scale. Finally, experimental validation using HCA tests on Toyoura sand is presented to corroborate the simulation findings.


% ================================================================
\section{DEM simulation setup}
\label{sec:ch4_dem_setup}

% ----------------------------------------------------------------
\subsection{DEM model of the hollow cylinder apparatus}
\label{subsec:ch4_dem_hca}

DEM simulations were performed in \textsc{PFC3D} \citep{Itasca2021} to reproduce undrained cyclic torsional shear in an HCA. Unlike periodic-cell element tests, the present HCA model uses rigid boundaries (inner and outer cylinders, top and bottom planes) and six torsional blades to apply shear while preserving laboratory-relevant boundary conditions.

The specimen geometry (Fig.~\ref{fig:ch4_specimen_generation}) follows a laboratory-scale HCA configuration with inner diameter \SI{6}{cm}, outer diameter \SI{10}{cm}, and specimen height \SI{10}{cm} \citep{Vargas2020}. A rolling-resistance contact model \citep{IwashitaOda1998,WensrichKatterfeld2012} is adopted to capture non-spherical grain effects, with a particle size distribution in the range of \SIrange{1.0}{3.0}{mm}, which is a 10$\times$ scaled-up analogue of the grain size distribution of Toyoura sand. The DEM parameters are summarized in Table~\ref{tab:ch4_dem_parameters}. Particles are generated by gravity deposition from the upper part of the apparatus. Gravity is then removed to ensure a uniform stress state, and six vertically arranged torsional blades are inserted into the specimen (Fig.~\ref{fig:ch4_specimen_generation}b). Compared with the conventional approach of constraining boundary particles to transmit torque \citep[e.g.,][]{Li2014,Liu2021}, these independent wall-based blades avoid interference with the radial expansion or contraction of the cylinders, thereby preserving accurate radial stress measurement.

\begin{table}[t]
  \centering
  \caption{Parameters in DEM simulation.}
  \label{tab:ch4_dem_parameters}
  \begin{tabular}{@{}lr@{}}
    \toprule
    Description & Value \\
    \midrule
    Number of particles, $N_p$                          & 53{,}764 \\
    Particle density, $\rho_s$ (\si{kg/m^3})            & 2650 \\
    Young's modulus, $E$ (\si{GPa})                     & 0.3 \\
    Normal-to-shear stiffness ratio, $\kappa$           & 2.0 \\
    Interparticle friction, $\mu_{pp}$                  & 0.0\textsuperscript{a}\,/\,0.50\textsuperscript{b} \\
    Particle--wall friction, $\mu_{pw}$                 & 0.0 \\
    Normal critical damping ratio, $\beta_n$            & 0.7 \\
    Shear critical damping ratio, $\beta_s$             & 0.5 \\
    Rolling friction coefficient, $\mu_r$               & 0.1 \\
    \bottomrule
    \multicolumn{2}{@{}l}{\footnotesize \textsuperscript{a}Isotropic consolidation stage.} \\
    \multicolumn{2}{@{}l}{\footnotesize \textsuperscript{b}Anisotropic consolidation and cyclic shear stages.}
  \end{tabular}
\end{table}

\begin{figure}[t]
  \centering
  \begin{tikzpicture}
    % --- (a) Deposit ---
    \node[anchor=north west, inner sep=0] (imgA) at (0,0)
      {\includegraphics[height=6cm]{specimen_generation_deposit.png}};
    \node[font=\small\bfseries, anchor=north west] at
      ([xshift=4pt,yshift=-4pt]imgA.north west) {(a)};
    % --- (b) Post-insert ---
    \node[anchor=north west, inner sep=0] (imgB) at ([xshift=8pt]imgA.north east)
      {\includegraphics[height=6cm]{specimen_generation_postinsert.png}};
    \node[font=\small\bfseries, anchor=north west] at
      ([xshift=4pt,yshift=-4pt]imgB.north west) {(b)};
    % --- Colorbar ---
    \node[anchor=north west, inner sep=0] (cbar) at ([xshift=8pt]imgB.north east)
      {
\includegraphics[height=6cm]{colorbar_viridis.png}};
    \foreach \val/\frac in {1.5/0, 1.25/0.25, 1.0/0.5, 0.75/0.75, 0.5/1.0} {
      \node[font=\footnotesize, anchor=west] at
        ([xshift=2pt, yshift={-\frac*6cm}]cbar.north east) {\val};
    }
    \node[font=\footnotesize, rotate=90, anchor=south] at
      ([xshift=40pt]cbar.east) {Particle radius (mm)};
  \end{tikzpicture}
  \caption{Initial specimen generation process in the DEM-based hollow cylinder apparatus model.
    Particles are colored by radius.
    (a)~Particle generation modeling air-pluviation.
    (b)~Insertion of torsional blades.}
  \label{fig:ch4_specimen_generation}
\end{figure}

% ----------------------------------------------------------------
\subsubsection{Stress and strain evaluation in the HCA}
\label{subsubsec:ch4_hca_equations}

Table~\ref{tab:ch4_hca_equations} summarizes the stress and strain evaluation expressions used in the DEM-based hollow cylinder apparatus model. In the following, $R$ and $r$ denote the current outer and inner radii, $H$ and $H_0$ are the current and initial specimen heights, $T$ is the torque, $\theta$ is the twist angle, $p_i^{\prime}$, $p_o^{\prime}$, and $p_z^{\prime}$ are the radial and vertical effective pressures measured from particle--wall contact forces on the inner cylinder, outer cylinder, and top/bottom plates, respectively, $u_i$ and $u_o$ are the inner and outer radial displacements, and $f_{pw}^{i}$, $f_{pw}^{o}$, $f_{pw}^{t}$, $f_{pw}^{b}$ are particle--wall contact forces on the inner cylinder, outer cylinder, top plate (including its attached blades), and bottom plate (including its attached blades), respectively. In the constant-volume DEM approximation, all stresses computed from contact forces are effective stresses. The values of $\sigma_\theta^{\prime}$ and $\sigma_r^{\prime}$ are derived from equilibrium relations, $\varepsilon_z$ and $\gamma_{z\theta}$ are based on strain compatibility, and the remaining stress and strain expressions follow the assumption of linear elasticity \citep{Hight1983}.

\begin{table*}[t]
  \centering
  \caption{Equations of stress and strain in the hollow cylinder apparatus.}
  \label{tab:ch4_hca_equations}
  \begin{tabular}{@{}lll@{}}
    \toprule
    Component & Stress & Strain \\
    \midrule
    Inner pressure &
    $\displaystyle p_i^{\prime} = \frac{\sum\!\sqrt{{f_{pw_x}^{i}}^2+{f_{pw_y}^{i}}^2}}{2\pi r H}$ &
    --- \\[12pt]
    Outer pressure &
    $\displaystyle p_o^{\prime} = \frac{\sum\!\sqrt{{f_{pw_x}^{o}}^2+{f_{pw_y}^{o}}^2}}{2\pi R H}$ &
    --- \\[12pt]
    Vertical pressure &
    $\displaystyle p_z^{\prime} = \frac{\sum\!|f_{pw_z}^{t}|+\sum\!|f_{pw_z}^{b}|}{2\pi(R^2-r^2)}$ &
    --- \\[12pt]
    Vertical &
    $\sigma_z^{\prime} = p_z^{\prime}$ &
    $\displaystyle\varepsilon_z = -\frac{H-H_0}{H_0}$ \\[12pt]
    Circumferential &
    $\displaystyle\sigma_\theta^{\prime} = \frac{p_o^{\prime} R - p_i^{\prime} r}{R-r}$ &
    $\displaystyle\varepsilon_\theta = -\frac{u_o+u_i}{R+r}$ \\[12pt]
    Radial &
    $\displaystyle\sigma_r^{\prime} = \frac{p_o^{\prime} R + p_i^{\prime} r}{R+r}$ &
    $\displaystyle\varepsilon_r = -\frac{u_o-u_i}{R-r}$ \\[12pt]
    Shear &
    $\displaystyle\tau_{z\theta} = \frac{3T}{2\pi(R^3-r^3)}$ &
    $\displaystyle\gamma_{z\theta} = \frac{2\theta(R^3-r^3)}{3H(R^2-r^2)}$ \\[12pt]
    Major principal &
    $\displaystyle\sigma_1^{\prime} = \frac{\sigma_z^{\prime}+\sigma_\theta^{\prime}}{2}+\sqrt{\left(\frac{\sigma_z^{\prime}-\sigma_\theta^{\prime}}{2}\right)^{\!2}+\tau_{z\theta}^2}$ &
    $\displaystyle\varepsilon_1 = \frac{\varepsilon_z+\varepsilon_\theta}{2}+\sqrt{\left(\frac{\varepsilon_z-\varepsilon_\theta}{2}\right)^{\!2}+\left(\frac{\gamma_{z\theta}}{2}\right)^{\!2}}$ \\[12pt]
    Intermediate &
    $\sigma_2^{\prime} = \sigma_r^{\prime}$ &
    $\varepsilon_2 = \varepsilon_r$ \\[12pt]
    Minor principal &
    $\displaystyle\sigma_3^{\prime} = \frac{\sigma_z^{\prime}+\sigma_\theta^{\prime}}{2}-\sqrt{\left(\frac{\sigma_z^{\prime}-\sigma_\theta^{\prime}}{2}\right)^{\!2}+\tau_{z\theta}^2}$ &
    $\displaystyle\varepsilon_3 = \frac{\varepsilon_z+\varepsilon_\theta}{2}-\sqrt{\left(\frac{\varepsilon_z-\varepsilon_\theta}{2}\right)^{\!2}+\left(\frac{\gamma_{z\theta}}{2}\right)^{\!2}}$ \\[12pt]
    Deviatoric &
    $\displaystyle q = \sqrt{\tfrac{1}{2}\!\left[(\sigma_1^{\prime}-\sigma_2^{\prime})^2+(\sigma_2^{\prime}-\sigma_3^{\prime})^2+(\sigma_3^{\prime}-\sigma_1^{\prime})^2\right]}$ &
    $\displaystyle\varepsilon_q = \sqrt{\tfrac{2}{9}\!\left[(\varepsilon_1-\varepsilon_2)^2+(\varepsilon_2-\varepsilon_3)^2+(\varepsilon_3-\varepsilon_1)^2\right]}$ \\[6pt]
    \bottomrule
  \end{tabular}
\end{table*}

% ----------------------------------------------------------------
\subsection{Specimen preparation and anisotropic consolidation protocol}
\label{subsec:ch4_preparation}

Specimens are prepared using an isotropic-consolidation to anisotropic-consolidation (IC--AC) protocol. First, the assembly is compressed under zero interparticle friction to a target void ratio, which is determined through the iterative calibration procedure described in \cref{subsec:ch4_calibration}: approximately 0.736 for the dense state ($D_r \approx 90\%$) and 0.742 for the medium-dense state ($D_r \approx 75\%$). The interparticle friction coefficient is then restored to 0.5, and the specimen is consolidated from an isotropic state (\Kzero{} = 1.0, $\pmean=\SI{10}{kPa}$) to the target mean effective stress ($\pmean=\SI{100}{kPa}$) along a linear anisotropic stress path.

During the AC stage, $\pmean$ is increased from \SI{10}{kPa} to \SI{100}{kPa} through ten equally spaced target stress states along a linear anisotropic stress path. The resulting stress paths and void-ratio evolution for the three primary cases are shown in \cref{fig:ch4_ac_path}. The markers in \cref{fig:ch4_ac_path}(a) correspond to the prescribed target states; the solid lines between them represent the actual stress evolution, which exhibits minor fluctuations that arise because different \Kzero{} states require different amounts of axial strain to reach the same $\pmean$ increment---compression-state specimens (\Kzero{} = 0.5) undergo larger axial strain per step, causing the axial stress to lag slightly behind the radial stress during each increment. These fluctuations vanish at each target state, confirming that the servo accurately satisfies the prescribed stress conditions. The final void ratio differs only slightly among the three states (on the order of $10^{-3}$), which allows the comparison to focus on fabric-related effects rather than void-ratio differences. \Cref{fig:ch4_ac_specimens} shows the three specimens after consolidation, with force-chain networks and particle arrangements displayed side by side; the visual contrast in force-chain directionality reflects the different \Kzero{} states.

\begin{figure}[t]
  \centering
  \begin{subfigure}[t]{0.48\linewidth}
    \centering
    \includegraphics[width=\linewidth]{anisotropic_consolidation_stress_path.png}
    \caption{$\pmean$--$q$ stress path.}
    \label{fig:ch4_ac_path_a}
  \end{subfigure}
  \hfill
  \begin{subfigure}[t]{0.48\linewidth}
    \centering
    \includegraphics[width=\linewidth]{anisotropic_consolidation_void_ratio.png}
    \caption{$\pmean$--$e$ evolution.}
    \label{fig:ch4_ac_path_b}
  \end{subfigure}
  \caption{Stress and void-ratio evolution during anisotropic consolidation for the three representative \Kzero{} states.}
  \label{fig:ch4_ac_path}
\end{figure}

\begin{figure}[t]
  \centering
  \begin{tikzpicture}
    % --- (a) K0=0.5 ---
    \node[anchor=north west, inner sep=0] (imgA) at (0,0)
      {\includegraphics[height=5.2cm]{specimen_k0_05_after_ac.png}};
    \node[font=\small, anchor=north] at (imgA.south) {(a) \Kzero{} = 0.5};
    % --- (b) K0=1.0 ---
    \node[anchor=north west, inner sep=0] (imgB) at ([xshift=3pt]imgA.north east)
      {\includegraphics[height=5.2cm]{specimen_k0_10_after_ac.png}};
    \node[font=\small, anchor=north] at (imgB.south) {(b) \Kzero{} = 1.0};
    % --- (c) K0=2.0 ---
    \node[anchor=north west, inner sep=0] (imgC) at ([xshift=3pt]imgB.north east)
      {\includegraphics[height=5.2cm]{specimen_k0_20_after_ac.png}};
    \node[font=\small, anchor=north] at (imgC.south) {(c) \Kzero{} = 2.0};
    % --- Colorbar ---
    \node[anchor=north west, inner sep=0] (cbar) at ([xshift=3pt]imgC.north east)
      {
\includegraphics[height=5.2cm]{colorbar_viridis_force.png}};
    \foreach \val/\frac in {2.0/0, 1.5/0.25, 1.0/0.5, 0.5/0.75, 0/1.0} {
      \node[font=\footnotesize, anchor=west] at
        ([xshift=2pt, yshift={-\frac*5.2cm}]cbar.north east) {\val};
    }
    \node[font=\footnotesize, rotate=90, anchor=south] at
      ([xshift=34pt]cbar.east) {Contact force magnitude (N)};
  \end{tikzpicture}
  \caption{DEM HCA specimens after anisotropic consolidation to $\pmean=\SI{100}{kPa}$.
    The left half of each panel shows interparticle force chains colored by contact force magnitude.
    (a)~\Kzero{} = 0.5; (b)~\Kzero{} = 1.0; (c)~\Kzero{} = 2.0.}
  \label{fig:ch4_ac_specimens}
\end{figure}

The cyclic shear simulation program consists of two sets. In the primary set, three representative \Kzero{} states---0.5 (compression), 1.0 (isotropic), and 2.0 (extension)---are examined at a dense state ($D_r \approx 90\%$) across multiple CSR levels to characterize how liquefaction resistance varies with the direction and magnitude of stress anisotropy, and to trace the associated energy dissipation, stiffness degradation, and fabric evolution. In the expanded set, a medium-dense state ($D_r \approx 75\%$) is introduced alongside the dense state, and two additional \Kzero{} values (0.67 and 1.5) are included to refine the coverage of anisotropy levels, yielding ten cases ($2 \times 5$) at a single CSR = 0.300. This expanded set allows examination of whether the \Kzero{}-dependent liquefaction resistance trend observed at $D_r \approx 90\%$ persists at a lower density, and provides a broader parametric basis for the subsequent microscopic analysis.

% ----------------------------------------------------------------
\subsection{Coupled analytical servo mechanism for undrained loading}
\label{subsec:ch4_servo}

In a laboratory HCA, the specimen is enclosed between an outer chamber (pressure $p_o$) and an inner chamber (pressure $p_i$), which are controlled independently. An axial load is applied through the loading piston, producing a vertical pressure $p_z$, and a torque $T$ is transmitted through end platens. Under undrained conditions the pore water pressure $u$, monitored by pressure transducers, is uniform throughout the specimen, so the effective pressures are $p_o^{\prime}=p_o-u$, $p_i^{\prime}=p_i-u$, and $p_z^{\prime}=p_z-u$. Because $u$ drops out in any difference of effective pressures, the radial pressure difference equals its total-stress counterpart, $p_o^{\prime}-p_i^{\prime}=p_o-p_i$, and so does the axial pressure difference between the vertical and average radial effective stresses (Table~\ref{tab:ch4_hca_equations}): $p_z^{\prime}-\sigma_r^{\prime}=p_z-\sigma_r$.

\begin{figure}[t]
  \centering
  \includegraphics[width=\linewidth]{fig_hca_standalone.pdf}
  \caption{Schematic of the hollow cylinder apparatus (HCA) specimen.
    The left half of each panel shows the laboratory boundary conditions with total stresses
    ($p_o$, $p_i$, $p_z$);
    the right half shows the DEM representation with effective stresses
    ($p_o^\prime$, $p_i^\prime$, $p_z^\prime$).
    (a)~Perspective view with specimen height $H$.
    (b)~Plan view showing the annular geometry with inner radius $r$,
    outer radius $R$, and torque $T$.}
  \label{fig:ch4_hca_schematic}
\end{figure}

The constant-volume method is a well-established approach in DEM for reproducing undrained cyclic loading \citep{YimsiriSoga2010}. In this method, the pore water pressure is not modelled explicitly; instead, the contact forces between particles and walls yield effective stresses ($p_o^{\prime}$, $p_i^{\prime}$, $p_z^{\prime}$) directly (Table~\ref{tab:ch4_hca_equations}). The DEM-based HCA specimen has four kinematic degrees of freedom---inner radius rate ($dr/dt$), outer radius rate ($dR/dt$), height rate ($dH/dt$), and rotation angle rate ($d\theta/dt$), where $r$ is the inner radius, $R$ the outer radius (Fig.~\ref{fig:ch4_hca_schematic}b), and $H$ the height of the hollow cylindrical specimen (Fig.~\ref{fig:ch4_hca_schematic}a)---and the following laboratory boundary conditions to reproduce: the undrained (constant-volume) constraint, the inner and outer cell pressures, the axial stress, and the torsional shear stress.

Among these, the constant-volume constraint provides one equation linking $dr/dt$, $dR/dt$, and $dH/dt$, while the torsional condition is satisfied by adjusting the rotation rate $d\theta/dt$ to match the target torque. The remaining challenge is to map the three effective pressures accessible in DEM---$p_i^{\prime}$, $p_o^{\prime}$, and $p_z^{\prime}$---to the laboratory total-stress boundary conditions, where the unknown pore water pressure $u$ prevents a direct one-to-one correspondence. The difference relations established above resolve this difficulty: because $u$ drops out, the two differences $p_o^{\prime}-p_i^{\prime}=p_o-p_i$ and $p_z^{\prime}-\sigma_r^{\prime}=p_z-\sigma_r$ map the DEM effective stresses to the laboratory total-stress conditions without requiring $u$, yielding two independent equations that, together with the volume constraint and the torsional condition, close the system of four equations for four unknowns.

\subsubsection{Boundary conditions}

\paragraph{Radial condition.}
The radial condition controls the pressure difference between the outer and inner cylinders. The radial pressure difference and its servo error are defined as:
\begin{equation}
p_{dif,r}^{\prime} = p_o^{\prime} - p_i^{\prime}, \qquad e_r = p_{dif,r}^{\prime} - p_{dif,r}^{tar},
\label{eq:ch4_cond1_radial}
\end{equation}
where $p_o^{\prime}$ and $p_i^{\prime}$ are the outer and inner radial effective pressures (Table~\ref{tab:ch4_hca_equations}) and $p_{dif,r}^{tar}$ is the target radial pressure difference. In this study, $p_o = p_i$ (equal chamber pressures) is adopted, so $p_{dif,r}^{tar} = 0$. The servo must drive $e_r$ to zero at each timestep.

\paragraph{Axial condition.}
Laboratory HCA tests can operate under either constant total axial stress or constant specimen height.
(i)~\emph{Stress-control mode}: the servo controls the difference between the vertical and radial effective stresses. The axial pressure difference and its servo error are:
\begin{equation}
p_{dif,z}^{\prime} = p_z^{\prime} - \sigma_r^{\prime}, \qquad e_z = p_{dif,z}^{\prime} - p_{dif,z}^{tar},
\label{eq:ch4_cond2_axial}
\end{equation}
where $p_z^{\prime}$ is the vertical effective pressure, $\sigma_r^{\prime}$ is the average radial effective stress, and $p_{dif,z}^{tar}$ is the target axial pressure difference. This mode is used during anisotropic consolidation and monotonic shear validation (\cref{subsec:ch4_calibration}).
(ii)~\emph{Displacement-control mode}: the specimen height is held constant ($dH/dt=0$), and $p_{dif,z}^{\prime}$ evolves freely. This mode is adopted for all cyclic shear simulations.

\paragraph{Torsional condition.}
The torque error is:
\begin{equation}
e_T = T - T^{tar},
\label{eq:ch4_cond3_torsion_error}
\end{equation}
where $T$ is the current torque and $T^{tar}$ is the target torque.

\paragraph{Volume constraint.}
The undrained (constant-volume) condition requires (with the outward-positive sign convention $dr>0$, $dR>0$ denoting radius increases):
\begin{equation}
2\pi H\!\left(R\frac{dR}{dt} - r\frac{dr}{dt}\right) + \pi\!\left(R^2-r^2\right)\frac{dH}{dt} = 0.
\label{eq:ch4_const_volume}
\end{equation}

\subsubsection{Analytical solution of the coupled servo}

The four conditions above---radial (\cref{eq:ch4_cond1_radial}), axial (\cref{eq:ch4_cond2_axial}), torsional (\cref{eq:ch4_cond3_torsion_error}), and constant-volume (\cref{eq:ch4_const_volume})---provide four independent equations for the four kinematic unknowns ($dr/dt$, $dR/dt$, $dH/dt$, $d\theta/dt$). The system is therefore closed and can be solved analytically. An analytical solution is obtained by reducing the four unknowns to two through two successive eliminations.

\paragraph{Step 1: Torsion (independent).}
The torsional condition does not enter the kinematic constraints that govern $dr/dt$, $dR/dt$, and $dH/dt$. It is therefore solved independently at each timestep:
\begin{equation}
\frac{d\theta}{dt} = \frac{e_T}{\Delta t}\,S_T,
\label{eq:ch4_torsion_servo}
\end{equation}
where $\Delta t$ is the simulation timestep and $S_T$ is the torsional servo coefficient, determined from the contact-stiffness-weighted rotational stiffness $I_{rot}$ about the specimen axis (\cref{fig:ch4_servo_geometry}):
\begin{equation}
I_{rot} = \sum k_n\bigl[(\mathbf{r}_d\times\hat{\mathbf{z}})\cdot\mathbf{n}\bigr]^2, \qquad S_T = \frac{1}{I_{rot}},
\label{eq:ch4_irot}
\end{equation}
where $\mathbf{r}_d$ is the position vector from the rotation axis to the contact point in the horizontal plane, $\hat{\mathbf{z}}$ is the unit vector along the specimen axis, $\mathbf{n}$ is the contact normal directed from the particle centre to the blade, and $k_n$ is the normal contact stiffness.

\begin{figure}[t]
  \centering
  \begin{tikzpicture}[
    every node/.style={font=\small},
    lbl/.style={font=\footnotesize},
  ]
    \node[anchor=south west, inner sep=0] (img) at (0,0)
      {\includegraphics[width=0.5\linewidth]{servo_torsional_contact_geometry.png}};
    \begin{scope}[
      x={(img.south east)},
      y={(img.north west)},
    ]
      \coordinate (CP) at (0.646, 0.276);
      \coordinate (RA) at (0.4, 0.80);
      \coordinate (RA_h) at (0.4, 0.56);

      \node[lbl, red!70!black, above right] at (0.46, 0.4) {$\mathbf{n}$};

      \draw[-{Stealth[length=3pt]}, thick, blue!70!black] (RA_h) -- (CP);
      \node[lbl, blue!70!black, above] at ($(RA_h)!0.52!(CP)$) {$\mathbf{r}_d$};
      \draw[-{Stealth[length=3pt]}, thick, black] (RA_h) -- (RA);
      \node[lbl, black, right] at ($(RA_h)!0.7!(RA)$) {$\hat{\mathbf{z}}$};
    \end{scope}
  \end{tikzpicture}
  \caption{Contact geometry used to evaluate the rotational stiffness contribution for torsional servo control. $\mathbf{r}_d$: position vector from the rotation axis to the contact point in the horizontal plane; $\mathbf{n}$: contact normal directed from particle centre to blade; $\hat{\mathbf{z}}$: unit vector along the specimen axis.}
  \label{fig:ch4_servo_geometry}
\end{figure}

\paragraph{Step 2: Volume constraint --- reduce to two unknowns.}
From the constant-volume constraint (\cref{eq:ch4_const_volume}), the outer radius rate is determined by:
\begin{equation}
\frac{dR}{dt} = \frac{r}{R}\frac{dr}{dt} - \frac{R^2-r^2}{2RH}\frac{dH}{dt}.
\label{eq:ch4_dRdt}
\end{equation}
After this elimination, the system reduces to two unknowns ($dr/dt$, $dH/dt$) with two target conditions ($e_r=0$ and $e_z=0$).

\subsubsection{Derivation of coupled servo coefficients}

Movement of both the inner wall ($dr$) and the top/bottom plates ($dH$) induces a change in the outer radius through \cref{eq:ch4_dRdt}, which in turn affects both $p_{dif,r}^{\prime}$ and $p_{dif,z}^{\prime}$. The servo coefficients must therefore account for this cross-coupling.

Let $K_{r,i}=\sum_{c\in\mathrm{inner}} k_n$ and $K_{r,o}=\sum_{c\in\mathrm{outer}} k_n$ denote the aggregate normal contact stiffness at the inner and outer cylinders, respectively, and $K_z=\frac{1}{2}\sum_{c\in\mathrm{top/bottom}} k_n$ the effective axial contact stiffness, where the summations run over all particle--wall contacts on the respective rigid boundaries. The factor $\frac{1}{2}$ arises because the top and bottom plates move symmetrically (each by $dH/2$) to achieve the total height change $dH$. The corresponding wall areas are $A_i=2\pi rH$, $A_o=2\pi RH$, and $A_z=2\pi(R^2-r^2)$.

The sign convention adopts outward-positive for all kinematic increments: $dr>0$ and $dR>0$ denote outward (radius-increasing) wall movements, and $dH>0$ denotes an increase in specimen height. The individual pressure increments are:
\begin{align}
\delta p_i^{\prime} &= \frac{K_{r,i}}{A_i}\,dr,\qquad
\delta p_o^{\prime} = -\frac{K_{r,o}}{A_o}\,dR, \label{eq:ch4_dp_r}\\[4pt]
\delta p_z^{\prime} &= -\frac{K_z}{A_z}\,dH, \label{eq:ch4_dp_z}
\end{align}
where the signs follow from the geometry: an outward movement of the inner wall ($dr>0$) expands the inner cylinder into the specimen, compressing particles against it and increasing $p_i^{\prime}$; an outward movement of the outer wall ($dR>0$) moves it away from the specimen, unloading particles and decreasing $p_o^{\prime}$; and an increase in specimen height ($dH>0$) moves the top and bottom plates apart, unloading axial contacts and decreasing $p_z^{\prime}$.

Substituting \cref{eq:ch4_dRdt} into the pressure increment expressions to eliminate $dR$ yields the pressure responses in terms of only $dr$ and $dH$. Collecting the coefficients into a matrix:
\begin{equation}
\begin{pmatrix} \delta p_{dif,r}^{\prime} \\[4pt] \delta p_{dif,z}^{\prime} \end{pmatrix}
=
\underbrace{\begin{pmatrix} \hat{K}_{11} & \hat{K}_{12} \\[8pt] \hat{K}_{21} & \hat{K}_{22} \end{pmatrix}}_{\hat{\mathbf{K}}}
\begin{pmatrix} dr \\[4pt] dH \end{pmatrix},
\label{eq:ch4_stiffness_matrix}
\end{equation}
where
\begin{alignat}{2}
\hat{K}_{11} &= -\frac{K_{r,i}}{A_i}-\frac{r}{R}\,\frac{K_{r,o}}{A_o}, &\qquad
\hat{K}_{12} &= \frac{R^2-r^2}{2RH}\,\frac{K_{r,o}}{A_o}, \label{eq:ch4_Khat_row1}\\[4pt]
\hat{K}_{21} &= -\frac{K_{r,i}-\frac{r}{R}\,K_{r,o}}{A_i+A_o}, &\qquad
\hat{K}_{22} &= -\frac{K_z}{A_z}-\frac{(R^2-r^2)\,K_{r,o}}{2RH\,(A_i+A_o)}. \label{eq:ch4_Khat_row2}
\end{alignat}
All entries carry dimensions of Pa/m, reflecting the displacement-to-pressure mapping through the wall areas. Three entries are negative ($\hat{K}_{11}$, $\hat{K}_{21}$, $\hat{K}_{22}<0$), consistent with the outward-positive convention: outward wall movements generally reduce the confining pressures. The exception is $\hat{K}_{12}>0$, because an increase in specimen height forces the outer wall inward through the volume constraint (\cref{eq:ch4_dRdt}), which raises $p_o^{\prime}$ and thereby increases $p_{dif,r}^{\prime}$.

The servo must apply wall movements that eliminate the current pressure errors, i.e., produce $\delta p_{dif,r}^{\prime}=-e_r$ and $\delta p_{dif,z}^{\prime}=-e_z$. Inverting \cref{eq:ch4_stiffness_matrix} gives the coupled servo equations:
\begin{align}
\frac{dr}{dt} &= S_{rr}\,\frac{e_r}{\Delta t} + S_{rz}\,\frac{e_z}{\Delta t}, \label{eq:ch4_servo_dr}\\[6pt]
\frac{dH}{dt} &= S_{zr}\,\frac{e_r}{\Delta t} + S_{zz}\,\frac{e_z}{\Delta t}, \label{eq:ch4_servo_dH}
\end{align}
where the servo coefficients are the entries of $\hat{\mathbf{K}}^{-1}$:
\begin{alignat}{2}
S_{rr} &= \frac{-\hat{K}_{22}}{\det(\hat{\mathbf{K}})}, &\qquad
S_{rz} &= \frac{\hat{K}_{12}}{\det(\hat{\mathbf{K}})}, \label{eq:ch4_servo_coeff_r}\\[4pt]
S_{zr} &= \frac{\hat{K}_{21}}{\det(\hat{\mathbf{K}})}, &\qquad
S_{zz} &= \frac{-\hat{K}_{11}}{\det(\hat{\mathbf{K}})}. \label{eq:ch4_servo_coeff_z}
\end{alignat}
with $\det(\hat{\mathbf{K}}) = \hat{K}_{11}\hat{K}_{22} - \hat{K}_{12}\hat{K}_{21}$. The diagonal terms $S_{rr}$ and $S_{zz}$ represent the direct servo response, while the off-diagonal terms $S_{rz}$ and $S_{zr}$ capture the cross-coupling between the radial and axial conditions. The remaining kinematic variable $dR/dt$ is then obtained by substituting $dr/dt$ and $dH/dt$ into \cref{eq:ch4_const_volume}. The exact analytical solution ensures that the servo coefficients are physically consistent, providing the correct basis on which a relaxation factor ($<1$) can be applied for numerical stability.

Under the constant specimen height condition ($dH/dt=0$) adopted for all cyclic shear simulations, only the radial condition remains active and the servo coefficient reduces to:
\begin{equation}
S_{rr} = -\frac{1}{\hat{K}_{11}} = \frac{1}{\dfrac{K_{r,i}}{A_i}+\dfrac{r}{R}\,\dfrac{K_{r,o}}{A_o}}.
\label{eq:ch4_Srr_constantH}
\end{equation}

The servo procedure applied at each simulation timestep is summarized in \cref{fig:ch4_servo_flowchart}.

\begin{figure}[t]
  \centering
  \begin{tikzpicture}[
    node distance=10pt,
    box/.style={draw, rounded corners=2pt, minimum width=0.88\linewidth,
                minimum height=7mm, align=center, font=\small},
    arr/.style={-{Stealth[length=5pt]}, thick},
    label/.style={font=\footnotesize\itshape, text=black!70},
  ]
    \node[box] (step1)
      {Compute $p_i^{\prime},\;p_o^{\prime},\;p_z^{\prime},\;T$
       from particle--wall contacts};
    \node[box, below=of step1] (step2)
      {Evaluate errors:\;
       $e_r=p_{dif,r}^{\prime}-p_{dif,r}^{tar}$,\;
       $e_z=p_{dif,z}^{\prime}-p_{dif,z}^{tar}$,\;
       $e_T=T-T^{tar}$};
    \node[box, below=of step2] (step3)
      {Torsion (independent):\;
       $d\theta/dt = (e_T/\Delta t)\,S_T$};
    \node[box, below=of step3] (step4)
      {Assemble $\hat{\mathbf{K}}$ from contact stiffnesses;\;
       compute $S_{rr},\,S_{rz},\,S_{zr},\,S_{zz}$};
    \node[box, below=of step4] (step5)
      {Coupled radial--axial servo:\;
       $dr/dt$,\; $dH/dt$ \;(\cref{eq:ch4_servo_dr,eq:ch4_servo_dH})};
    \node[box, below=of step5] (step6)
      {Volume constraint $\rightarrow$ $dR/dt$
       \;(\cref{eq:ch4_const_volume})};
    \node[box, below=of step6] (step7)
      {Move walls;\; advance timestep};

    \draw[arr] (step1) -- (step2);
    \draw[arr] (step2) -- (step3);
    \draw[arr] (step3) -- (step4);
    \draw[arr] (step4) -- (step5);
    \draw[arr] (step5) -- (step6);
    \draw[arr] (step6) -- (step7);
    \draw[arr] (step7.south) -- ++(0,-0.3) -| ([xshift=8pt]step1.east |- step7.south) -- ([xshift=8pt]step1.east) -- (step1.east)
      node[label, pos=0.45, right=20pt, rotate=90, anchor=south] {next timestep};
  \end{tikzpicture}
  \caption{Servo procedure at each simulation timestep during undrained loading.}
  \label{fig:ch4_servo_flowchart}
\end{figure}

\paragraph{Excess pore water pressure.}
The excess pore water pressure is inferred from the reduction in radial effective stress:
\begin{equation}
u = \sigma_{r0}^{\prime} - \sigma_r^{\prime},
\label{eq:ch4_epwp}
\end{equation}
where $\sigma_r^{\prime}$ is the current radial effective stress and $\sigma_{r0}^{\prime}$ is its initial value. The ratio $\ru$ is then defined as $\ru = u/\sigma_{r0}^{\prime}$.

\paragraph{Comparison with earlier formulations.}
Compared with earlier DEM-based HCA servo strategies, the present formulation makes three contributions explicit. First, the governing equations are systematically reduced through elimination of the volume constraint and torsional independence before the coupled system is solved, clarifying the mathematical structure of the problem. Second, the laboratory-to-DEM mapping is established through two pressure differences ($p_{dif,r}^{\prime}$ and $p_{dif,z}^{\prime}$), which are the correct servo targets because they are invariant to the unknown pore water pressure. Third, the cross-coupling between the radial and axial degrees of freedom is quantified analytically through the off-diagonal entries of $\hat{\mathbf{K}}$, so that the servo coefficients correctly reflect the mutual influence of $dr/dt$ and $dH/dt$ on both boundary conditions. None of these steps were formalized in earlier work: \citet{Han2024jgs} solved only the radial and torsional conditions while keeping the specimen height fixed, and \citet{Ma2024} introduced an axial servo but derived each servo coefficient from the contact stiffness at the corresponding wall independently, without reflecting the cross-coupling among the wall movements.

% ----------------------------------------------------------------
\subsection{Parameter calibration and validation}
\label{subsec:ch4_calibration}

The micro-parameters listed in Table~\ref{tab:ch4_dem_parameters} were calibrated against the undrained monotonic torsional shear data of \citet{Nakata1998} for dense Toyoura sand ($D_r \approx 90\%$, $\pmean_0=\SI{100}{kPa}$, \Kzero{} = 1.0) \citep[see also][]{Ma2024}. Because the micro-parameters collectively influence shear stiffness, dilatancy, and liquefaction resistance, the calibration targeted these behaviors in sequence. Starting from a base parameter set at a given void ratio, the contact Young's modulus $E$ and rolling friction coefficient $\mu_r$ were adjusted while keeping the normal-to-shear stiffness ratio $\kappa$ constant, aiming to reproduce the shear stiffness observed in the $\varepsilon_q$--$q$ response of the monotonic shear tests. The effective stress path in $\pmean$--$q$ space was then examined to assess dilatancy. If the simulated dilatancy was insufficient---i.e., the mean effective stress at phase transformation was lower than the experimental value---the specimen was further compacted to a lower void ratio, which increased both dilatancy and shear stiffness; $E$ was then reduced accordingly to restore the target shear stiffness. This iterative process between void ratio, contact modulus, and rolling friction was repeated until both the shear stiffness and dilatancy simultaneously matched the monotonic experimental data. During the monotonic calibration, the axial pressure difference was maintained at zero ($p_{dif,z}^{tar}=0$ in \cref{eq:ch4_cond2_axial}) to ensure $\sigma_z^{\prime}=\sigma_r^{\prime}$ during shearing.

The calibrated parameter set was then validated against the cyclic liquefaction resistance data of \citet{Ishihara1985} for the same sand under identical conditions. Specimens were subjected to sinusoidal torsional shear at five CSR levels (0.20, 0.25, 0.30, 0.35, and 0.40) under constant specimen height ($dH/dt=0$), matching the experimental protocol, with liquefaction defined by single-amplitude shear strain $\gamma_{SA}=2.5\%$.

\begin{figure}[t]
  \centering
  \includegraphics[width=0.95\linewidth]{validate_combined.png}
  \caption{Calibration and validation of the DEM-HCA model for dense Toyoura sand ($D_r \approx 90\%$, $\pmean_0=\SI{100}{kPa}$, \Kzero{} = 1.0): (a)~effective stress path and (b)~stress--strain relationship from undrained monotonic shear, calibrated against \citet{Nakata1998}; (c)~cyclic stress ratio versus number of cycles to liquefaction, validated against \citet{Ishihara1985}.}
  \label{fig:ch4_validation}
\end{figure}

\Cref{fig:ch4_validation} summarizes the calibration and validation results. The effective stress path (\cref{fig:ch4_validation}a) captures the initial contraction followed by dilation towards the critical state line, matching the experimental trend. The stress--strain relationship (\cref{fig:ch4_validation}b) reproduces the shear stiffness and strain-hardening behavior. The liquefaction resistance curve (\cref{fig:ch4_validation}c) shows good agreement with the cyclic data across the full range of CSR. The simultaneous agreement across stiffness, dilatancy, and cyclic liquefaction resistance confirms the reliability of the micro-parameters and demonstrates that the servo mechanism performs well under both axial stress-control (monotonic calibration) and displacement-control (cyclic validation) modes.

% ----------------------------------------------------------------
\subsection{Microscopic descriptors}
\label{subsec:ch4_micro_descriptors}

While the macroscopic response characterizes the overall liquefaction behavior, a key advantage of DEM is its ability to access particle-scale information that is not directly measurable in laboratory tests. Three microscopic descriptors are introduced to probe the contact network from complementary perspectives.

\paragraph{Mechanical coordination number.}
The mechanical coordination number \citep{Thornton2000,Oda1977} is defined as:
\begin{equation}
Z_m = \frac{2N_c - N_1}{N_p - N_1 - N_0},
\label{eq:ch4_zm}
\end{equation}
where $N_c$ is the total number of interparticle contacts, $N_p$ is the total number of particles, and $N_1$ and $N_0$ are the numbers of particles with one and zero contacts, respectively. Floaters are excluded to focus on the load-bearing contact network. $Z_m$ serves as a scalar measure of packing compactness.

\paragraph{Fabric tensor and fabric anisotropy indicator.}
The directionality of the contact network is characterized through the fabric tensor \citep{Oda1982fabric,Oda1982stat}:
\begin{equation}
\Phi_{ij} = \frac{1}{N_c}\sum n_i\,n_j,
\label{eq:ch4_fabric_tensor}
\end{equation}
where $n_i$ and $n_j$ are the components of the contact normal vector. Because the HCA specimen is axisymmetric about the vertical axis $z$, it is natural to express $\Phi_{ij}$ in the cylindrical coordinate system $(r,\theta,z)$.

The deviatoric fabric tensor $\mathbf{F}_{ij}$ and its second invariant $\Fc$ measure the magnitude of fabric anisotropy \citep{ZhaoGuo2013}:
\begin{equation}
\mathbf{F}_{ij} = \frac{15(\Phi_{ij} - \delta_{ij}/3)}{2}, \qquad
\Fc = \sqrt{\frac{3}{2}\mathbf{F}_{ij}\mathbf{F}_{ij}}.
\label{eq:ch4_Fc}
\end{equation}

However, $\Fc$ captures only the magnitude of anisotropy without indicating its direction. A signed indicator is needed to distinguish between compression and extension states. Following \citet{Otsubo2022}, the fabric anisotropy indicator $\alpha$ \citep{YimsiriSoga2010} is defined from the axial diagonal component $\Phi_{zz}$:
\begin{equation}
\alpha = \frac{5(3\Phi_{zz} - 1)}{5\Phi_{zz} + 1}.
\label{eq:ch4_alpha}
\end{equation}
Positive $\alpha$ indicates that contact normals concentrate toward the vertical axis (compression-state anisotropy, \Kzero{} $< 1.0$), whereas negative $\alpha$ indicates concentration toward the horizontal plane (extension-state anisotropy, \Kzero{} $> 1.0$).

\paragraph{Contact density.}
The contact density \citep{Han2023} describes the average number of contacts per unit solid angle on the unit sphere of contact normal orientations:
\begin{equation}
\rho_c(\theta_z,\,\phi_{\mathrm{cir}}) = \frac{2\,N(\theta_z,\,\phi_{\mathrm{cir}})}{N_p\,\Delta\Omega},
\label{eq:ch4_contact_density}
\end{equation}
where $\theta_z$ and $\phi_{\mathrm{cir}}$ are the polar and azimuthal angles in the spherical coordinate system, $N(\theta_z,\,\phi_{\mathrm{cir}})$ is the number of contacts whose normals fall within the angular bin of widths $\Delta\theta_z$ and $\Delta\phi_{\mathrm{cir}}$, and $\Delta\Omega = (\cos\phi_{\mathrm{cir}} - \cos(\phi_{\mathrm{cir}}+\Delta\phi_{\mathrm{cir}}))\,\Delta\theta_z$ is the solid angle subtended by that bin. Unlike the fabric tensor, $\rho_c$ retains the full directional information and thus provides a more complete picture of the evolving contact network.

Together, these three descriptors probe the contact network at increasing levels of directional detail: \Zm{} quantifies the average contact density per particle (scalar), $\alpha$ captures the dominant anisotropy direction (signed scalar), and $\rho_c$ resolves the full directional distribution (field on the unit sphere).


% ================================================================
\section{Results and discussion: DEM simulation}
\label{sec:ch4_results_sim}

% ----------------------------------------------------------------
\subsection{Macroscopic response and liquefaction resistance}
\label{subsec:ch4_macro}

The macroscopic response is examined first to establish the overall liquefaction resistance ordering among \Kzero{} states. \Cref{fig:ch4_macro_response} presents a combined overview of the typical macroscopic response during undrained cyclic torsional shear at CSR = 0.200.

\begin{figure}[t]
  \centering
  \begin{subfigure}[t]{0.48\linewidth}
    \centering
    \includegraphics[width=\linewidth]{macro_response_a.png}
    \caption{$\tau_{z\theta}$ vs.\ $\pmean$ for \Kzero{} = 1.0.}
    \label{fig:ch4_macro_response_a}
  \end{subfigure}
  \hfill
  \begin{subfigure}[t]{0.48\linewidth}
    \centering
    \includegraphics[width=\linewidth]{macro_response_b.png}
    \caption{$\tau_{z\theta}$ vs.\ $\gamma_{z\theta}$ for \Kzero{} = 1.0.}
    \label{fig:ch4_macro_response_b}
  \end{subfigure}

  \vspace{0.5em}

  \begin{subfigure}[t]{0.48\linewidth}
    \centering
    \includegraphics[width=\linewidth]{macro_response_c.png}
    \caption{Deviatoric stress vs.\ $\pmean$ for three \Kzero{} states.}
    \label{fig:ch4_macro_response_c}
  \end{subfigure}
  \hfill
  \begin{subfigure}[t]{0.48\linewidth}
    \centering
    \includegraphics[width=\linewidth]{macro_response_d.png}
    \caption{Excess pore pressure ratio $\ru$ evolution for three \Kzero{} states.}
    \label{fig:ch4_macro_response_d}
  \end{subfigure}
  \caption{Typical macroscopic response during undrained cyclic torsional shear at CSR = 0.200.}
  \label{fig:ch4_macro_response}
\end{figure}

As shown in \cref{fig:ch4_macro_response}(a,\,b), the isotropically consolidated specimen (\Kzero{} = 1.0) exhibits the classical undrained cyclic response: progressive reduction in $\pmean$ with a butterfly-shaped effective stress path, and marked stiffness degradation with increasingly nonlinear hysteretic loops. As the shear stress $\tau_{z\theta}$ cyclically acts on the hollow cylindrical specimen, the mean effective stress $\pmean$ exhibits an overall decrease. Once $\pmean$ falls below approximately \SI{60}{kPa}, an increase in the magnitude of $\tau_{z\theta}$ drives an upward trend in $\pmean$, displaying the characteristic butterfly-shaped relationship. After liquefaction onset, $\pmean$ and $\tau_{z\theta}$ delineate the critical state line and converge at the origin periodically. The stress--strain relationship (\cref{fig:ch4_macro_response}b) exhibits stiff shear modulus with slight reduction in the initial stage; as cyclic loading progresses towards liquefaction, shear stiffness markedly decreases, plastic deformation occurs, and strong nonlinearity becomes evident in the hysteretic loops.

When the three \Kzero{} states are compared in \cref{fig:ch4_macro_response}(c,\,d), their stress paths converge toward a similar low-effective-stress state near the origin, but the rate of effective stress loss differs systematically: the \Kzero{} = 2.0 specimen accumulates excess pore pressure fastest, the \Kzero{} = 0.5 specimen slowest, and the \Kzero{} = 1.0 specimen is intermediate. The \Kzero{} = 1.0 specimen starts with deviatoric stress $q$ at zero due to isotropic consolidation, while \Kzero{} = 0.5 and \Kzero{} = 2.0 specimens exhibit initial deviatoric stress due to anisotropic consolidation. As cyclic shear progresses, specimens with different \Kzero{} values follow distinct stress paths toward the critical state line, with varying rates of stress degradation. Initial stress anisotropy therefore governs the rate of effective stress reduction rather than the ultimate state at liquefaction.

Two commonly used liquefaction criteria are considered: a stress-based criterion, defined as excess pore water pressure reaching 95\% of the initial radial effective stress $\sigma_{r0}^{\prime}$ (i.e., $\ru = 0.95$), and a strain-based criterion, defined as single-amplitude shear strain reaching $\gamma_{SA} = 2.5\%$ (as adopted in \cref{subsec:ch4_calibration} for validation against \citet{Ishihara1985}). A comparison of the two criteria across all simulated cases confirms that they yield nearly identical conclusions: the strain-based criterion consistently produces $N_L$ values approximately 1--2 cycles higher than the stress-based criterion (e.g., for CSR = 0.200, $N_L$ = 40.7 vs.\ 39.5 for \Kzero{} = 0.5, 32.8 vs.\ 31.5 for \Kzero{} = 1.0, and 30.7 vs.\ 29.5 for \Kzero{} = 2.0), and the relative ordering among \Kzero{} states remains unchanged. Since this study focuses on comparing specimens with different initial stress states, the stress-based criterion ($\ru = 0.95$) is adopted for all subsequent analyses, as it directly reflects the loss of effective stress that constitutes the fundamental mechanism of liquefaction.

\begin{figure}[t]
  \centering
  \includegraphics[width=0.7\linewidth]{liquefaction_resistance_curves.png}
  \caption{Cyclic liquefaction resistance curves for specimens with three representative stress anisotropies (\Kzero{} = 0.5, 1.0, 2.0). The inset shows the expanded \Kzero{} set at CSR = 0.300.}
  \label{fig:ch4_liq_resistance}
\end{figure}

\Cref{fig:ch4_liq_resistance} presents the CSR--$N_L$ relationship for the three representative anisotropic states. Across all CSR levels, $N_L$ decreases monotonically from \Kzero{} = 0.5 through 1.0 to 2.0; the inset at CSR = 0.300, which includes two additional intermediate states (\Kzero{} = 0.67 and 1.5), confirms this trend. The effect of \Kzero{} depends on the direction of deviatoric stress: compression-type anisotropy (\Kzero{} $<$ 1.0) enhances liquefaction resistance, while extension-type anisotropy (\Kzero{} $>$ 1.0) reduces it. This directional dependency is consistent with the findings by \citet{GeorgiannouKonstadinou2014} for dense sand states and the experimental results of \citet{Vargas2020}, who reported approximately 20\% higher liquefaction strength for \Kzero{} = 0.5 specimens compared to isotropic conditions. More recently, \citet{Zhao2024} conducted systematic HCA torsional shear tests covering both compressional and extensional consolidation states and reported higher shear strengths under compressive anisotropic consolidation, further supporting this directional dependency. This contrasts with some earlier experimental studies \citep{IshiharaTakatsu1979,YamashitaToki1993} that reported minimal \Kzero{} dependency. While these studies established the macroscopic trend, the particle-scale origin of this directional dependency is examined in \cref{subsec:ch4_micro}.

% ----------------------------------------------------------------
\subsection{Energy and stiffness interpretations}
\label{subsec:ch4_energy_stiffness}

To interpret the macroscopic origin of the observed \Kzero{}-dependent liquefaction resistance, two complementary perspectives are considered: cumulative shear work (energy input) and shear-wave-velocity evolution (stiffness degradation).

\subsubsection{Cumulative shear work}

Cumulative shear work refers to the energy input during undrained cyclic shear. It is valuable to examine the correlation between the liquefaction resistance of soils and their susceptibility to the input energy. \citet{TowhataIshihara1985} conducted a series of experiments where isotropically consolidated sand specimens were subjected to undrained cyclic shear under various combinations of torsional and axial shear stresses with continuous rotation of principal stress axes. They revealed a unique relationship between excess pore water pressure and shear work that is independent of the shear stress paths applied. \citet{Figueroa1994} similarly confirmed that the shear work required for triggering liquefaction is independent of the amplitude of strain through strain-controlled tests. \citet{GeorgiannouKonstadinou2014} concluded from the comparison between IC and AC specimens that the energy associated with terminal water pressure is positively correlated with relative density; for equivalent relative densities, AC specimens require greater energy than IC specimens to induce liquefaction.

The cumulative unit volume shear work $W_s$ is defined as \citep{TowhataIshihara1985}:
\begin{equation}
W_s = \int_0^t \dot{\gamma}_{z\theta}\,\tau_{z\theta}\,dt,
\label{eq:ch4_shear_work}
\end{equation}
where $\dot{\gamma}_{z\theta}$ is the shear strain rate and $\tau_{z\theta}$ is the shear stress. Note that $W_s$ represents the total shear work input computed from boundary quantities, rather than the plastic (dissipated) component alone; however, because the recoverable elastic strain energy becomes negligibly small relative to the cumulative work as loading progresses, the two measures follow closely similar trends throughout the cyclic loading history.

\begin{figure}[t]
  \centering
  \begin{subfigure}[t]{0.48\linewidth}
    \centering
    \includegraphics[width=\linewidth]{energy_ru_vs_normalized_cycles.png}
    \caption{$\ru$ vs.\ normalized cycle ratio $N/N_L$.}
    \label{fig:ch4_energy_work_a}
  \end{subfigure}
  \hfill
  \begin{subfigure}[t]{0.48\linewidth}
    \centering
    \includegraphics[width=\linewidth]{cumulative_work_vs_normalized_cycles.png}
    \caption{$W_s$ vs.\ normalized cycle ratio $N/N_L$.}
    \label{fig:ch4_energy_work_b}
  \end{subfigure}

  \vspace{0.5em}

  \begin{subfigure}[t]{0.48\linewidth}
    \centering
    \includegraphics[width=\linewidth]{cumulative_work_vs_cycles.png}
    \caption{$W_s$ vs.\ actual number of cycles $N$.}
    \label{fig:ch4_energy_work_c}
  \end{subfigure}
  \hfill
  \begin{subfigure}[t]{0.48\linewidth}
    \centering
    \includegraphics[width=\linewidth]{ru_vs_cumulative_work.png}
    \caption{$\ru$ vs.\ cumulative shear work $W_s$.}
    \label{fig:ch4_energy_work_d}
  \end{subfigure}
  \caption{Excess pore pressure ratio and cumulative shear work evolution during cyclic shear, and their relationship for different \Kzero{} states (CSR = 0.200).}
  \label{fig:ch4_energy_work}
\end{figure}

\Cref{fig:ch4_energy_work} summarizes the evolution of $\ru$ and $W_s$ for the three \Kzero{} states at CSR = 0.200. In normalized coordinates (\cref{fig:ch4_energy_work}(a,\,b)), both $\ru$ and $W_s$ follow similar growth patterns across the three \Kzero{} states and converge as liquefaction is approached, indicating that the energy accumulation process follows a common progression toward liquefaction. In the pre-liquefaction stage, the three states separate in $\ru$ (\cref{fig:ch4_energy_work}a): the \Kzero{} = 0.5 specimen consistently exhibits the lowest $\ru$, the \Kzero{} = 2.0 specimen the highest, and the \Kzero{} = 1.0 specimen falls in between. In contrast, the normalized $W_s$ curves (\cref{fig:ch4_energy_work}b) show a different ordering: the \Kzero{} = 0.5 specimen accumulates slightly more energy per normalized cycle while the \Kzero{} = 1.0 specimen accumulates the least.

This apparent discrepancy arises because the normalization by $N_L$ compresses different numbers of actual cycles onto the same horizontal axis. When $W_s$ is plotted against the actual cycle count $N$ (\cref{fig:ch4_energy_work}c), the ordering aligns with the liquefaction resistance trend: the \Kzero{} = 2.0 specimen accumulates energy fastest, \Kzero{} = 1.0 is intermediate, and \Kzero{} = 0.5 accumulates energy most slowly. This indicates that the key difference among anisotropic states lies in how rapidly energy is input per cycle, rather than in the total energy required to trigger liquefaction.

Having examined $\ru$ and $W_s$ separately, a natural next step is to test whether their direct relationship can explain the \Kzero{}-dependent liquefaction resistance, following the energy-based approach of \citet{TowhataIshihara1985}, who showed a unique $\ru$--$W_s$ relationship for isotropically consolidated specimens under different shear stress paths. \Cref{fig:ch4_energy_work}(d) shows that the \Kzero{} = 0.5 specimen exhibits a distinctly lower $\ru$ at any given $W_s$ compared to the other two states, which partially explains its higher liquefaction resistance---the compression-state fabric requires more energy input to generate the same level of pore pressure ratio. However, when comparing \Kzero{} = 1.0 and \Kzero{} = 2.0, no clear distinction emerges in the $\ru$--$W_s$ relationship despite their different liquefaction resistance (\cref{fig:ch4_liq_resistance}), limiting the explanatory power of this energy-based approach. It should be noted that the comparison has inherent limitations: \citet{TowhataIshihara1985} demonstrated a unique $W_s$--$u$ relationship for specimens under the same isotropic initial consolidation state but subjected to different shear stress paths, whereas the present study compares specimens with different initial stress anisotropies under the same loading path. Moreover, the quantity compared here is $\ru$ (EPWP normalized by the initial lateral stress) rather than the absolute pore water pressure $u$, and the different initial stress states result in different normalization bases. These factors suggest that the energy-based interpretation alone is insufficient to fully account for the \Kzero{} effect on liquefaction resistance. Adopting the strain-based criterion ($\gamma_{SA}=2.5\%$) in place of $\ru=0.95$ does not alter this conclusion, as both criteria yield the same $N_L$ ordering among \Kzero{} states.

\subsubsection{Shear-wave velocity evolution and stiffness degradation}

The inability of the energy-based approach to distinguish between \Kzero{} = 1.0 and \Kzero{} = 2.0 suggests that the difference in liquefaction resistance between these two states arises not from how much energy is needed to generate pore pressure, but from how rapidly energy is input per cycle. Since the per-cycle energy input is governed by strain amplitude, which in turn depends on shear stiffness, this motivates an examination of the relationship between shear stiffness evolution and liquefaction resistance. The shear-wave velocity \Vs{} is examined as a per-cycle stiffness proxy. Unlike cumulative energy, \Vs{} is evaluated within each quarter-cycle and therefore tracks the current stiffness state of the specimen as loading progresses, providing a direct link between the macroscopic energy accumulation rate and the underlying particle-scale fabric evolution discussed in \cref{subsec:ch4_micro}. The shear-wave velocity is evaluated from the secant shear modulus $G$ and saturated density $\rho_{\mathrm{sat}}$ as \citep{XuLing2015}:
\begin{equation}
V_s = \sqrt{\frac{G}{\rho_{\mathrm{sat}}}},
\label{eq:ch4_vs}
\end{equation}
where $G$ is the secant shear modulus and $\rho_{\mathrm{sat}}$ is the saturated density.

\begin{figure}[t]
  \centering
  \includegraphics[width=0.7\linewidth]{shear_wave_velocity_evolution.png}
  \caption{Evolution of shear-wave velocity during cyclic shear for different \Kzero{} values, with loading and unloading stiffness shown separately (CSR = 0.200).}
  \label{fig:ch4_vs_evolution}
\end{figure}

The loading $V_s$ is calculated from the secant modulus of the first and third quarter-periods of each cycle, and the unloading $V_s$ from the second and fourth quarter-periods. \Cref{fig:ch4_vs_evolution} shows that unloading stiffness is consistently higher than loading stiffness for all three states, reflecting the typical hysteretic behavior of granular materials under cyclic loading. The \Kzero{} = 1.0 specimen exhibits the highest initial \Vs{}, with values of approximately \SI{138}{m/s} (loading) and \SI{148}{m/s} (unloading). The \Kzero{} = 0.5 and \Kzero{} = 2.0 specimens show lower and comparable initial values, with loading \Vs{} around \SIrange{125}{127}{m/s} and unloading \Vs{} around \SIrange{142}{144}{m/s}. However, the subsequent stiffness degradation differs markedly among the three states. During the early stage ($N < 15$), the \Kzero{} = 1.0 specimen maintains the highest stiffness. As cycling progresses, the degradation rates diverge: the \Kzero{} = 1.0 and \Kzero{} = 2.0 specimens degrade relatively rapidly, while the \Kzero{} = 0.5 specimen degrades much more slowly. By approximately 20 cycles, the \Kzero{} = 0.5 specimen surpasses \Kzero{} = 1.0 in loading \Vs{}, and by 28 cycles, the contrast becomes striking---the \Kzero{} = 0.5 specimen retains a loading \Vs{} of approximately \SI{103}{m/s}, compared to \SI{80}{m/s} for \Kzero{} = 1.0 and \SI{51}{m/s} for \Kzero{} = 2.0.

These different degradation trajectories directly explain the per-cycle energy accumulation rates observed in \cref{fig:ch4_energy_work}(c): the \Kzero{} = 0.5 specimen, whose stiffness degrades most slowly, develops the smallest strain amplitude per cycle and thus accumulates energy most gradually; the \Kzero{} = 2.0 specimen, with the fastest stiffness degradation, develops the largest strain amplitude per cycle and accumulates energy most rapidly; and the \Kzero{} = 1.0 specimen exhibits intermediate behavior. Compared with the energy-based approach, which cannot distinguish between \Kzero{} = 1.0 and 2.0 (\cref{fig:ch4_energy_work}d), the stiffness degradation rate provides a unified macroscopic explanation for the liquefaction resistance ordering across all three \Kzero{} states. Notably, the initial stiffness alone cannot predict these outcomes---the \Kzero{} = 0.5 and \Kzero{} = 2.0 specimens share comparable initial \Vs{} yet exhibit markedly different degradation trajectories. It is the rate of stiffness degradation during cyclic loading, rather than the initial stiffness, that ultimately governs energy accumulation behavior and liquefaction resistance.

While stiffness degradation provides a consistent macroscopic explanation, it does not reveal why different \Kzero{} states exhibit different degradation rates. These macroscopic observations point to fundamental differences in the underlying particle-scale fabric structure. To elucidate the microscopic mechanisms, the following sections examine the evolution of coordination number, fabric anisotropy, and contact density under different \Kzero{} states.

% ----------------------------------------------------------------
\subsection{Microscopic fabric interpretation}
\label{subsec:ch4_micro}

The preceding macroscopic analysis identified stiffness degradation rate as the proximate cause of \Kzero{}-dependent liquefaction resistance. This subsection examines the underlying particle-scale fabric.

\subsubsection{Mechanical coordination number}

\begin{figure}[t]
  \centering
  \includegraphics[width=0.7\linewidth]{coordination_number_evolution.png}
  \caption{Mechanical coordination number evolution during cyclic shear and relationship between initial \Zm{} and the number of cycles to liquefaction (CSR = 0.200).}
  \label{fig:ch4_zm_evolution}
\end{figure}

\Cref{fig:ch4_zm_evolution} shows that the initial coordination number differs with stress anisotropy: the isotropically consolidated specimen exhibits the largest $Z_{m0} \approx 4.98$, whereas the anisotropically consolidated specimens have slightly smaller and nearly identical values ($Z_{m0} \approx 4.89$ and $4.88$ for \Kzero{} = 0.5 and 2.0, respectively). This is consistent with the larger initial stiffness of the isotropic state. During cyclic loading, \Zm{} decreases gradually in the early stage and then drops more rapidly as liquefaction is approached, reflecting progressive loss of contacts. Liquefaction is triggered when \Zm{} approaches a similar threshold value ($\sim$4.0) across the three states, and post-liquefaction \Zm{} fluctuates around comparable levels.

However, the inset in \cref{fig:ch4_zm_evolution} also indicates that \Zm{} alone cannot explain the \Kzero{} dependence of liquefaction resistance: the \Kzero{} = 0.5 and 2.0 specimens exhibit similar initial \Zm{} but markedly different $N_L$, implying that directional characteristics of the contact network must be considered.

\subsubsection{Fabric anisotropy indicator}

\begin{figure}[t]
  \centering
  \includegraphics[width=0.7\linewidth]{fabric_anisotropy_evolution.png}
  \caption{Evolution of fabric anisotropy indicator $\alpha$ during cyclic shear for different \Kzero{} states (CSR = 0.200).}
  \label{fig:ch4_alpha_evolution}
\end{figure}

\Cref{fig:ch4_alpha_evolution} shows that the \Kzero{} = 0.5 specimen starts with positive $\alpha_0 \approx +0.07$, indicating contact normals concentrated in the axial direction. The \Kzero{} = 1.0 specimen shows $\alpha_0$ close to zero ($\approx +0.02$), reflecting a nearly isotropic fabric. The \Kzero{} = 2.0 specimen exhibits negative $\alpha_0 \approx -0.05$, indicating contact normals concentrated in the horizontal direction. During the early stages of cyclic loading ($N/N_L < 0.8$), $\alpha$ remains relatively stable for all three \Kzero{} states. As the specimens approach liquefaction, $\alpha$ begins to fluctuate significantly and converge toward a narrow band (approximately 0 to 0.05), indicating that the directional fabric anisotropy diminishes as the fabric structure degrades. Notably, the ordering of initial values ($\alpha_0$: \Kzero{} = 0.5 $>$ 1.0 $>$ 2.0) largely mirrors the liquefaction resistance ranking, suggesting that initial fabric anisotropy is a key factor controlling liquefaction resistance. This relationship is examined quantitatively with an expanded dataset in \cref{subsec:ch4_combined_effect}.

While the fabric anisotropy indicator $\alpha$ provides a scalar measure of directional bias, the morphological basis for the sign of $\alpha_0$ can be visualized through contact density distributions. To simultaneously capture changes in both direction and quantity of contact normals during cyclic shear, the contact density $\rho_c$ (\cref{eq:ch4_contact_density}) is used for visualization.

\subsubsection{Contact density evolution and morphological interpretation}

\begin{figure}[t]
  \centering
  \begin{tikzpicture}
    \def\imgW{0.22\linewidth}
    \def\gap{1pt}
    \def\rowgap{12pt}

    % --- Row 1: K0=0.5 ---
    \node[anchor=north west, inner sep=0] (a1) at (0,0)
      {\includegraphics[width=\imgW]{contact_density_k0_05_t0.png}};
    \node[font=\small, anchor=south] at (a1.north) {$N/N_L = 0.00$};
    \node[anchor=north west, inner sep=0] (a2) at ([xshift=\gap]a1.north east)
      {\includegraphics[width=\imgW]{contact_density_k0_05_t302.png}};
    \node[font=\small, anchor=south] at (a2.north) {$N/N_L \approx 0.61$};
    \node[anchor=north west, inner sep=0] (a3) at ([xshift=\gap]a2.north east)
      {\includegraphics[width=\imgW]{contact_density_k0_05_t523.png}};
    \node[font=\small, anchor=south] at (a3.north) {$N/N_L \approx 1.05$};
    \node[anchor=north west, inner sep=0] (a4) at ([xshift=\gap]a3.north east)
      {\includegraphics[width=\imgW]{contact_density_k0_05_t529.png}};
    \node[font=\small, anchor=south] at (a4.north) {$N/N_L \approx 1.07$};

    % Axis labels
    \coordinate (axO) at ([xshift=6pt, yshift=12pt]a1.south west);
    \draw[->, thick] (axO) -- ++(0, 14pt) node[font=\scriptsize, right, inner sep=1pt] {$z$};
    \draw[->, thick] (axO) -- ++(16pt, 2pt) node[font=\scriptsize, right, inner sep=1pt] {$\theta$};
    \draw[->, thick] (axO) -- ++(8pt, -8pt) node[font=\scriptsize, right, inner sep=1pt] {$r$};

    % K0=0.5 label
    \path (a1.south) -- (a4.south) coordinate[midway] (row1mid);
    \node[font=\scriptsize, anchor=north] at ([yshift=-1pt]row1mid) {$\Kzero{} = 0.5$};

    % --- Row 2: K0=2.0 ---
    \node[anchor=north west, inner sep=0] (b1) at ([yshift=-\rowgap]a1.south west)
      {\includegraphics[width=\imgW]{contact_density_k0_20_t0.png}};
    \node[anchor=north west, inner sep=0] (b2) at ([xshift=\gap]b1.north east)
      {\includegraphics[width=\imgW]{contact_density_k0_20_t226.png}};
    \node[anchor=north west, inner sep=0] (b3) at ([xshift=\gap]b2.north east)
      {\includegraphics[width=\imgW]{contact_density_k0_20_t379.png}};
    \node[anchor=north west, inner sep=0] (b4) at ([xshift=\gap]b3.north east)
      {\includegraphics[width=\imgW]{contact_density_k0_20_t386.png}};

    % K0=2.0 label
    \path (b1.south) -- (b4.south) coordinate[midway] (row2mid);
    \node[font=\scriptsize, anchor=north] at ([yshift=-1pt]row2mid) {$\Kzero{} = 2.0$};

    % --- Rainbow colorbar ---
    \pgfpointdiff{\pgfpointanchor{b4}{south}}{\pgfpointanchor{a4}{north}}
    \pgfgetlastxy{\discard}{\cbarH}
    \node[anchor=north west, inner sep=0] (cbar) at ([xshift=4pt]a4.north east)
      {\includegraphics[height=\cbarH]{colorbar_rainbow.png}};
    \foreach \val/\frac in {0.5/0, 0.4/0.2, 0.3/0.4, 0.2/0.6, 0.1/0.8, 0/1.0} {
      \node[font=\footnotesize, anchor=west] at
        ([xshift=2pt, yshift={-\frac*\cbarH}]cbar.north east) {\val};
    }
    \node[font=\footnotesize, rotate=90, anchor=south] at
      ([xshift=32pt]cbar.east) {Contact density $\rho_c$};
  \end{tikzpicture}
  \caption{Evolution of contact density distribution during the liquefaction process for \Kzero{} = 0.5 and 2.0 specimens (CSR = 0.200).}
  \label{fig:ch4_contact_density}
\end{figure}

\Cref{fig:ch4_contact_density} shows that the \Kzero{} = 0.5 specimen initially exhibits an axially extended contact-density distribution with higher density (orange/red) at the poles and lower density (green) around the equator, corresponding to the positive $\alpha_0$ value. In contrast, the \Kzero{} = 2.0 specimen displays a dimpled sphere distribution with lower contact density at the poles and higher density around the equator, corresponding to the negative $\alpha_0$ value. These distinct morphologies arise from the stress-induced fabric rearrangement during anisotropic consolidation: contacts tend to align perpendicular to the maximum principal stress direction.

As cyclic shear progresses, the contact density distributions gradually evolve. During the intermediate stage, the overall contact density decreases notably while the morphological characteristics of the initial distribution are largely preserved---the fabric shrinks but does not reshape. As the specimens approach liquefaction, the contact density further diminishes and the direction of maximum contact density varies following the rotating principal stress axis. In the post-liquefaction state, the distinct initial morphologies converge toward similar patterns---both specimens exhibit an inclined distribution aligned with the rotated principal stress direction, rather than the original axial or horizontal orientations. This shared post-liquefaction regime is consistent with the dynamic convergence of $\alpha$, \Vs{}, and \Zm{} observed as specimens approach liquefaction.

\subsubsection{Interparticle force chains and particle displacement}

\begin{figure}[t]
  \centering
  \small
  \begin{tikzpicture}
    \def\imgW{0.188\linewidth}
    \def\gap{1pt}
    \def\rowgap{12pt}

    % --- Row 1: K0=0.5 ---
    \node[anchor=north west, inner sep=0] (a1) at (0,0)
      {\includegraphics[width=\imgW]{force_disp_k0_05_t0.png}};
    \node[font=\small, anchor=south] at (a1.north) {$N/N_L = 0.00$};
    \node[anchor=north west, inner sep=0] (a2) at ([xshift=\gap]a1.north east)
      {\includegraphics[width=\imgW]{force_disp_k0_05_t302.png}};
    \node[font=\small, anchor=south] at (a2.north) {$N/N_L \approx 0.61$};
    \node[anchor=north west, inner sep=0] (a3) at ([xshift=\gap]a2.north east)
      {\includegraphics[width=\imgW]{force_disp_k0_05_t523.png}};
    \node[font=\small, anchor=south] at (a3.north) {$N/N_L \approx 1.05$};
    \node[anchor=north west, inner sep=0] (a4) at ([xshift=\gap]a3.north east)
      {\includegraphics[width=\imgW]{force_disp_k0_05_t529.png}};
    \node[font=\small, anchor=south] at (a4.north) {$N/N_L \approx 1.07$};

    % K0=0.5 label
    \path (a1.south) -- (a4.south) coordinate[midway] (row1mid);
    \node[font=\scriptsize, anchor=north] at ([yshift=-1pt]row1mid) {$\Kzero{} = 0.5$};

    % --- Row 2: K0=2.0 ---
    \node[anchor=north west, inner sep=0] (b1) at ([yshift=-\rowgap]a1.south west)
      {\includegraphics[width=\imgW]{force_disp_k0_20_t0.png}};
    \node[anchor=north west, inner sep=0] (b2) at ([xshift=\gap]b1.north east)
      {\includegraphics[width=\imgW]{force_disp_k0_20_t226.png}};
    \node[anchor=north west, inner sep=0] (b3) at ([xshift=\gap]b2.north east)
      {\includegraphics[width=\imgW]{force_disp_k0_20_t379.png}};
    \node[anchor=north west, inner sep=0] (b4) at ([xshift=\gap]b3.north east)
      {\includegraphics[width=\imgW]{force_disp_k0_20_t386.png}};

    % K0=2.0 label
    \path (b1.south) -- (b4.south) coordinate[midway] (row2mid);
    \node[font=\scriptsize, anchor=north] at ([yshift=-1pt]row2mid) {$\Kzero{} = 2.0$};

    % --- Compute full span height for colorbars ---
    \pgfpointdiff{\pgfpointanchor{b1}{south}}{\pgfpointanchor{a1}{north}}
    \pgfgetlastxy{\discard}{\cbarH}

    % --- Force colorbar (left, spanning both rows) ---
    \node[anchor=north east, inner sep=0] (fcbar) at ([xshift=-4pt]a1.north west)
      {
\includegraphics[height=\cbarH]{colorbar_viridis_force.png}};
    \foreach \val/\frac in {2.0/0, 1.5/0.25, 1.0/0.5, 0.5/0.75, 0/1.0} {
      \node[font=\footnotesize, anchor=east] at
        ([xshift=-2pt, yshift={-\frac*\cbarH}]fcbar.north west) {\val};
    }
    \node[font=\footnotesize, rotate=90, anchor=south] at
      ([xshift=-18pt]fcbar.west) {Contact force (N)};

    % --- Displacement colorbar (right, spanning both rows) ---
    \node[anchor=north west, inner sep=0] (dcbar) at ([xshift=4pt]a4.north east)
      {
\includegraphics[height=\cbarH]{colorbar_sand_disp.png}};
    \foreach \val/\frac in {1.0/0, 0.75/0.25, 0.5/0.5, 0.25/0.75, 0/1.0} {
      \node[font=\footnotesize, anchor=west] at
        ([xshift=2pt, yshift={-\frac*\cbarH}]dcbar.north east) {\val};
    }
    \node[font=\footnotesize, rotate=90, anchor=south] at
      ([xshift=28pt]dcbar.east) {Displacement (mm)};
  \end{tikzpicture}
  \caption{Evolution of force chains (left) and particle displacements (right) during liquefaction for \Kzero{} = 0.5 and 2.0 specimens (CSR = 0.200).}
  \label{fig:ch4_force_disp}
\end{figure}

\Cref{fig:ch4_force_disp} visualizes the complementary evolution of contact force chains and particle displacements. At the initial state, force chains exhibit distinct patterns corresponding to the initial stress anisotropy. For \Kzero{} = 0.5, the contact forces are predominantly aligned in the axial direction, while for \Kzero{} = 2.0, they are concentrated in the horizontal direction, reflecting the maximum principal stress direction. As cyclic loading progresses, the number of interparticle contacts gradually decreases and the average contact force magnitude diminishes. The amplitude of particle displacement also increases with each loading cycle, reflecting the progressive degradation of the fabric structure.

As the specimens approach liquefaction, the force chains become increasingly sparse while the contact forces reorient to align with the rotated principal stress direction. In the post-liquefaction state, regardless of the initial \Kzero{} value, the contact force distributions converge toward similar patterns---forces no longer align with the initial axial or horizontal directions but concentrate along the rotated principal stress axis. Particle displacements likewise converge in magnitude and pattern, with larger displacements concentrated near the top and bottom boundaries of the specimen. These visual observations directly corroborate the statistical evidence from the contact density and fabric anisotropy analyses: cyclic shearing progressively erases the memory of the initial stress state, driving both assemblies into a shared post-liquefaction dynamic equilibrium. This convergence of fabric structure provides microscopic evidence supporting the uniqueness of the critical state stress ratio, although it should be noted that the post-liquefaction state observed here does not constitute the critical state in a strict sense, as the critical state requires sustained large deformation at the void ratio and effective stress corresponding to the critical state line.

% ----------------------------------------------------------------
\subsection{Combined effect of contact density and directional anisotropy on liquefaction resistance}
\label{subsec:ch4_combined_effect}

The qualitative analyses above show that initial \Zm{} alone cannot explain the \Kzero{}-dependent liquefaction resistance, while $\alpha$ provides additional discriminating power. To separate their respective contributions quantitatively, the expanded simulation set (two relative densities, five \Kzero{} values, CSR = 0.300) is examined.

\begin{figure}[t]
  \centering
  \begin{subfigure}[t]{0.48\linewidth}
    \centering
    \includegraphics[width=\linewidth]{fabric_liquefaction_2d_zm_nl.png}
    \caption{$Z_{m0}$ vs.\ $N_L$.}
    \label{fig:ch4_fabric_liq_zm}
  \end{subfigure}
  \hfill
  \begin{subfigure}[t]{0.48\linewidth}
    \centering
    \includegraphics[width=\linewidth]{fabric_liquefaction_2d_alpha_nl.png}
    \caption{$\alpha_0$ vs.\ $N_L$.}
    \label{fig:ch4_fabric_liq_alpha}
  \end{subfigure}
  \caption{Relationship between initial microscopic fabric parameters and liquefaction resistance $N_L$ for five \Kzero{} states at two relative densities (CSR = 0.300). Labels indicate \Kzero{} values.}
  \label{fig:ch4_fabric_liq_2d}
\end{figure}

\begin{figure}[t]
  \centering
  \includegraphics[width=0.7\linewidth]{fabric_liquefaction_3d.png}
  \caption{Three-dimensional view of the combined effect of $Z_{m0}$ and $\alpha_0$ on $N_L$ (CSR = 0.300).}
  \label{fig:ch4_fabric_liq_3d}
\end{figure}

The most immediate observation from \cref{fig:ch4_fabric_liq_2d}(a) is that $D_r \approx 90\%$ specimens consistently exhibit higher $N_L$ than $D_r \approx 75\%$ specimens across all five \Kzero{} states. Correspondingly, $Z_{m0}$ is systematically higher for $D_r \approx 90\%$ specimens (4.89--4.98) than for $D_r \approx 75\%$ specimens (4.77--4.88). This is consistent with the well-established understanding that denser packings develop more inter-particle contacts and thus greater liquefaction resistance. However, this cross-density correlation should be attributed to the coupling between macroscopic compactness and $Z_{m0}$ rather than an independent microscopic effect. Within each density group, $Z_{m0}$ peaks near \Kzero{} = 1.0 and decreases toward both compression and extension states, yet does not correlate monotonically with $N_L$. For instance, among $D_r \approx 90\%$ specimens, \Kzero{} = 0.5 exhibits the highest $N_L$ but not the highest $Z_{m0}$, while \Kzero{} = 2.0 has a comparable $Z_{m0}$ to \Kzero{} = 0.5 yet demonstrates considerably lower $N_L$. This confirms that $N_L$ is not independently controlled by $Z_{m0}$.

By contrast, \cref{fig:ch4_fabric_liq_2d}(b) shows that $\alpha_0$ correlates positively with $N_L$ within both density groups: compression-state anisotropy (larger $\alpha_0$) is associated with higher liquefaction resistance, while extension-state anisotropy (smaller $\alpha_0$) is associated with lower resistance. A minor reversal at \Kzero{} = 1.5 and 2.0 in the lower-density group, where $N_L$ differs by only one cycle, does not affect the overall trend. The two density groups follow a similar pattern along the $\alpha_0$ axis, indicating that the influence of fabric anisotropy on liquefaction resistance is consistent regardless of packing density.

The three-dimensional view in \cref{fig:ch4_fabric_liq_3d} synthesizes these observations. The ten data points form two distinct clusters separated primarily along the $Z_{m0}$ axis, while within each cluster $N_L$ varies systematically with $\alpha_0$. Both microscopic parameters contribute to liquefaction resistance, albeit along different dimensions: $Z_{m0}$ characterizes the overall contact density related to packing state, accounting for the inter-group separation between the two relative densities, whereas $\alpha_0$ captures the directional bias of the contact network induced by anisotropic consolidation, accounting for the intra-group variation among different \Kzero{} states. Notably, the two groups overlap in $Z_{m0}$: the $D_r = 75\%$, \Kzero{} = 1.0 specimen ($Z_{m0} \approx 4.88$, $\alpha_0 \approx +0.02$, $N_L \approx 4.5$) and the $D_r = 90\%$, \Kzero{} = 2.0 specimen ($Z_{m0} \approx 4.89$, $\alpha_0 \approx -0.05$, $N_L \approx 8.5$) share nearly identical $Z_{m0}$, yet the denser specimen resists liquefaction nearly twice as long despite its less favorable $\alpha_0$. This cross-group comparison indicates that relative density influences liquefaction resistance through additional mechanisms beyond what $Z_{m0}$ and $\alpha_0$ can capture---such as particle interlocking, force chain stability, and the distribution of contact forces---that collectively render the denser packing more resistant to cyclic loading regardless of the directional fabric bias. Neither parameter alone is sufficient; together, they provide a two-parameter microscopic characterization in which $Z_{m0}$ reflects packing compactness and $\alpha_0$ reflects the directional concentration of torsion-resisting contacts.

\subsubsection{Fabric tensor decomposition in cylindrical coordinates}

The physical origin of the $\alpha_0$--$N_L$ correlation can be understood by decomposing the fabric tensor in cylindrical coordinates (\cref{fig:ch4_fabric_mechanism}). Since $\tau_{z\theta}$ acts on the $z$- and $\theta$-planes, a contact's contribution to torsional resistance depends on the projection of its normal onto the $z$--$\theta$ plane; purely radially oriented contacts have zero projection and therefore do not contribute. The effective fraction of contacts resisting $\tau_{z\theta}$ is thus captured by $\Phi_{zz} + \Phi_{\theta\theta} = 1 - \Phi_{rr}$.

\begin{figure}[t]
  \centering
  \includegraphics[width=0.95\linewidth]{fabric_mechanism_schematic.png}
  \caption{Mechanism of torsion-resisting contact fraction under \Kzero{} consolidation. (a)~Schematic of a hollow cylindrical element in cylindrical coordinates: contacts with normals in the $z$ and $\theta$ directions (blue) directly resist the applied torsional shear $\tau_{z\theta}$, while radially oriented contacts (red) are mechanically neutral. (b)~Torsion-resisting fraction $1-\Phi_{rr}$ from the expanded dataset (Table~\ref{tab:ch4_fabric_diagonal}), showing a monotonic decrease with increasing \Kzero{} in both density groups.}
  \label{fig:ch4_fabric_mechanism}
\end{figure}

\begin{table}[t]
\centering
\caption{Initial fabric tensor diagonal components and torsion-resisting fraction for the expanded dataset (CSR = 0.300).}
\label{tab:ch4_fabric_diagonal}
\small
\begin{tabular}{cccccc}
\toprule
$D_r$ & \Kzero{} & $\Phi_{rr}$ & $\Phi_{\theta\theta}$ & $\Phi_{zz}$ & $1-\Phi_{rr}$ \\
\midrule
\multirow{5}{*}{90\%}
& 0.5  & 0.333 & 0.321 & 0.346 & 0.667 \\
& 0.67 & 0.337 & 0.323 & 0.340 & 0.664 \\
& 1.0  & 0.338 & 0.325 & 0.338 & 0.662 \\
& 1.5  & 0.339 & 0.326 & 0.335 & 0.661 \\
& 2.0  & 0.343 & 0.331 & 0.325 & 0.657 \\
\midrule
\multirow{5}{*}{75\%}
& 0.5  & 0.333 & 0.319 & 0.348 & 0.667 \\
& 0.67 & 0.338 & 0.322 & 0.340 & 0.662 \\
& 1.0  & 0.339 & 0.323 & 0.337 & 0.661 \\
& 1.5  & 0.341 & 0.325 & 0.334 & 0.659 \\
& 2.0  & 0.346 & 0.332 & 0.322 & 0.654 \\
\bottomrule
\end{tabular}
\end{table}

This fraction varies monotonically with \Kzero{} due to the redistribution of contacts under anisotropic consolidation. When \Kzero{} increases above unity, contacts migrate from the $z$ direction into both the $r$ and $\theta$ directions; the gain in $\theta$-directed contacts contributes to torsional resistance, but the accompanying gain in $r$-directed contacts does not. Conversely, when \Kzero{} decreases below unity, contacts concentrate axially, increasing the combined torsion-resisting fraction $\Phi_{zz} + \Phi_{\theta\theta}$. Extraction of the diagonal components from the expanded dataset (Table~\ref{tab:ch4_fabric_diagonal}) confirms that $\Phi_{rr}$ increases monotonically with \Kzero{} in both density groups, so that $1 - \Phi_{rr}$ decreases monotonically---in direct agreement with the liquefaction resistance ordering.

% ----------------------------------------------------------------
\subsection{Implications and limitations of DEM simulation}
\label{subsec:ch4_implications_sim}

The preceding analyses yield two broader implications.

First, the fabric tensor decomposition assumes that, in the HCA torsional shear geometry, contacts with normals in the $z$ and $\theta$ directions contribute equally to resistance against $\tau_{z\theta}$, so that only their sum $\Phi_{zz} + \Phi_{\theta\theta}$ matters and the individual partitioning between $\Phi_{zz}$ and $\Phi_{\theta\theta}$ is not considered. This assumption leads to a testable prediction: specimens sharing the same $\Phi_{zz} + \Phi_{\theta\theta}$ but differing in their individual partitioning should exhibit similar liquefaction resistance. Verification would require preparation methods that independently control the orientation of fabric anisotropy, which is left for future investigation.

It should also be noted that the foregoing interpretation is based on the geometric distribution of contact normals; however, anisotropic consolidation simultaneously induces contact force anisotropy, concentrating larger normal forces on contacts aligned with the major principal stress direction. Because these force-bearing contacts coincide with the torsion-resisting orientations for compression-state specimens, force anisotropy may amplify---or even dominate---the geometric fabric effect identified above. Quantitatively separating the contributions of contact-normal geometry and contact-force magnitude to the overall torsional resistance is left for future work.

One limitation should be noted. Relative density and $Z_{m0}$ are strongly coupled under the present consolidation procedure, making it difficult to vary one independently of the other. The present dataset therefore cannot fully isolate their individual contributions to liquefaction resistance; developing preparation protocols that decouple these variables remains a direction for future work.


% ================================================================
\section{Experimental validation}
\label{sec:ch4_experiment}

% TODO: Convert content from 04_2_anisotropic_consolidation_exp.md
% This section presents HCA experiments on Toyoura sand to validate
% the DEM simulation findings regarding K0 effects on liquefaction.

\subsection{Experimental apparatus}
\label{subsec:ch4_apparatus}

% TODO: Experimental setup description (pressure control, water injection,
%       torsional shear apparatus, servo system, data acquisition)

\subsection{Testing procedure}
\label{subsec:ch4_procedure}

% TODO: Specimen preparation, saturation, consolidation, cyclic loading protocol

\subsection{Experimental results}
\label{subsec:ch4_exp_results}

% TODO: Experimental liquefaction resistance curves, comparison with simulation

\subsection{Comparison between simulation and experiment}
\label{subsec:ch4_comparison}

% TODO: Systematic comparison of DEM and experimental findings


% ================================================================
\section{Summary}
\label{sec:ch4_summary}

This chapter used DEM simulations of undrained cyclic torsional shear in hollow cylindrical specimens to investigate how initial stress anisotropy (\Kzero{}) affects liquefaction resistance under HCA-type boundary conditions. The main findings are:

\begin{enumerate}
\item A coupled servo formulation was developed for the DEM-HCA model. The pressure differences ($p_{dif,r}'$ and $p_{dif,z}'$) were identified as the correct servo targets for mapping laboratory boundary conditions, and the constant-volume constraint was incorporated to couple the radial and axial degrees of freedom. Analytical inversion of the resulting stiffness matrix yields closed-form servo coefficients that ensure numerical stability.

\item Liquefaction resistance exhibits a monotonic directional dependency on initial stress anisotropy: compression-state anisotropy (\Kzero{} $< 1.0$) produces higher resistance, while extension-state anisotropy (\Kzero{} $> 1.0$) reduces it. This ordering is consistent across all CSR levels and both relative densities examined. The per-cycle energy input rate differs because of distinct stiffness degradation trajectories: the compression-state specimen degrades stiffness most slowly, and the extension-state specimen fastest.

\item The mechanical coordination number \Zm{} and the signed fabric anisotropy indicator $\alpha$ play complementary roles. \Zm{} primarily reflects packing compactness and separates the two density groups, but does not independently explain the \Kzero{}-dependent resistance within each group. By contrast, $\alpha_0$ correlates positively with $N_L$ within both density groups, indicating that directional fabric governs the resistance variation among \Kzero{} states. Decomposition of the fabric tensor in cylindrical coordinates shows that a contact's contribution to torsional resistance depends on the projection of its normal onto the $z$--$\theta$ plane; compression-state consolidation concentrates a larger fraction of contacts onto this plane, offering a geometric interpretation of the $\alpha_0$--$N_L$ correlation.

\item Multiple indicators, including shear wave velocity, mechanical coordination number, fabric anisotropy indicator, contact-density distributions, and force-chain configurations, all show that initially distinct fabric states progressively degrade and converge toward similar configurations in the post-liquefaction stage. This convergence indicates that cyclic loading drives the granular assembly into a common dynamic equilibrium regardless of the initial stress anisotropy.
\end{enumerate}

\end{document}
